\input _template

\setlanguage{EN}

\Title{Angle Chasing Basics}

\MathcompsLink{angle-basics-1}

\Author{Patrik Bak}

\sec Introduction

Geometry is the most visual part of mathematics -- we can draw it on paper, in software like \Link[https://www.geogebra.org/]{GeoGebra}, and we can enjoy seeing with our own eyes that things in it work. Why this is so, however, can be a mystery. To uncover these secrets, we need to build up a foundation of knowledge and visual intuition piece by piece.

One of the most fundamental techniques is to understand the world of angles and to know how to use their properties successfully in various types of problems. Only a few interesting examples don't use angles. The theory around angles is not complicated, but it is important to build it up slowly and in detail. We will start with the simplest properties, which don't even require working with circles. Even these can already be used successfully in many problems.

In this handout we will focus on angles, but we will deliberately leave out angles in circles, which deserve separate attention. For a moment, let's forget that circles exist.

To close the introduction, a small note about convention. In this handout we will see many diagrams of triangles in which vertex~$A$ is drawn at the top. It turns out that this way of drawing is the international standard, especially in the world of the Mathematical Olympiad and similar competitions. In Czechia and Slovakia, however, schools still customarily draw~$C$ at the top. In some countries one can also see~$B$ at the top, e.g. in Ukraine. From the point of view of solving problems, it is always a good idea to put at the top the vertex that makes the diagram look as symmetric as possible -- it is then easier to spot further symmetric things in it. In this handout the problems will be posed so that the symmetry is more visible with~$A$ at the top \SlightSmile.

\sec Basics of the World of Angles

In this section we will derive the simplest properties related to angles. I firmly believe that the key to mastering geometry is understanding things from the ground up and in depth, so we will spend more time on these simple properties.

Basic properties of angles, which we commonly use almost without thinking when solving problems:

\Highlight{
    \begitems \style i
    \i Vertical angles are equal:
       \Image{angles-vertical.pdf}
    \i Corresponding angles are equal:
       \Image{angles-corresponding.pdf}
    \i Alternate angles are equal:
       \Image{angles-alternate.pdf}
    \i Angles in a linear pair sum to $180^\circ$:
       \Image{angles-supplementary.pdf}
   \enditems
}

These properties are very naturally related to one another; a small exercise to ponder:

\Exercise{}{
    Convince yourself that:
    \begitems \style a
    \i the linear-pair property is equivalent to the vertical-angles property
    \i any two of the three properties about vertical, corresponding, and alternate angles imply the third
    \enditems
}{
    \textbf{(a)} The angle $\alpha'$ in a linear pair with a given angle $\alpha$ has measure $180^\circ - \alpha$. The angle in a linear pair with $\alpha'$ in turn has measure $180^\circ - (180^\circ - \alpha) = \alpha$. But this last angle is also vertical to $\alpha$, so vertical angles are equal.

    \Image{angles-supplementary-derive-vertical.pdf}

    Conversely, assume the vertical-angles property and label the consecutive angles at the intersection of two lines $\alpha$, $\beta$, $\gamma$, $\delta$. By the equality of vertical angles, $\alpha = \gamma$ and $\beta = \delta$. The sum of all four is a full angle, i.e. $\alpha + \beta + \gamma + \delta = 360^\circ$. Substituting, we get $2\alpha + 2\beta = 360^\circ$, so $\alpha + \beta = 180^\circ$ -- which is exactly the linear-pair property.

    \Image{angles-vertical-derive-supplementary.pdf}

    \textbf{(b)} All three properties say that a certain pair of angles has the same measure.

    Look at the three angles $\alpha$, $\beta$, $\gamma$ in the figure and at what each property says about their relationships. For the pair $\alpha$, $\beta$ at the upper intersection, the vertical-angles property gives $\alpha = \beta$. For the pair $\alpha$, $\gamma$, the corresponding-angles property gives $\alpha = \gamma$. Finally, for the pair $\beta$, $\gamma$, the alternate-angles property gives $\beta = \gamma$.

    \Image{angles-corresponding-alternate-supplementary-connection.pdf}

    The three properties therefore assert three equalities among the same three angles, just via different pairs. As soon as any two of them hold, the third follows by transitivity: e.g. if (V) and (C) hold, then $\alpha = \beta$ and $\alpha = \gamma$, so $\beta = \gamma$, which is (A). The other pairs are analogous.
}

The following statement is familiar to all of us, but can you prove it?

\Theorem{}{
    The sum of the angles in a triangle is $180^\circ$.
}{
    Through vertex~$A$ draw a line parallel to $BC$. Since $AB$ is a transversal of the parallel lines, the angle $\beta$ at~$A$ and the angle $\beta$ at~$B$ are alternate, hence equal. Similarly, via the transversal $AC$, the angles $\gamma$ at~$A$ and at~$C$ are alternate. The angles $\beta$, $\alpha$, $\gamma$ lie next to each other along the line through~$A$, so $\alpha + \beta + \gamma = 180^\circ.$

    \Image{angles-triangle.pdf}
}

This statement can certainly be generalized. We know, for example, that the sum of the angles in a quadrilateral is $360^\circ$. What about a pentagon? And a 67-gon? The answer is the following statement:

\Theorem{}{
    Let $n \ge 3$ be a natural number. The sum of the angles in a convex $n$-gon equals $(n-2) \cdot 180^\circ$.
}{
    \NamedProof{Proof 1 (mathematical induction).} Base case $n=3$: the sum of the angles of a triangle is $180^\circ = (3-2)\cdot 180^\circ$. Inductive step: let $A_1 A_2 \ldots A_{n+1}$ be a convex $(n+1)$-gon. The diagonal $A_1 A_n$ splits it into a triangle $A_1 A_n A_{n+1}$ and a convex $n$-gon $A_1 A_2 \ldots A_n$: the angles at vertices $A_2, \ldots, A_{n-1}$ belong entirely to the $n$-gon, the angle at $A_{n+1}$ belongs entirely to the triangle, and the angles at $A_1$, $A_n$ are split between the two figures so that their parts add up to the original interior angles of the $(n+1)$-gon. The angle sum of the $(n+1)$-gon is therefore
    $$
    \underbrace{180^\circ}_{\text{triangle}} + \underbrace{(n-2) \cdot 180^\circ}_{n\text{-gon}} = (n-1) \cdot 180^\circ = \bigl((n+1) - 2\bigr) \cdot 180^\circ.
    $$

    \Image{angles-polygon.pdf}

    \NamedProof{Proof 2 (fan triangulation).} From vertex $A_1$ draw diagonals to every non-adjacent vertex $A_3, A_4, \ldots, A_{n-1}$. This produces $n-2$ triangles that cover the interior of the $n$-gon without overlap. At each vertex $A_k$, the angles of the adjacent triangles at $A_k$ assemble exactly into the interior angle of the $n$-gon at $A_k$, so the sum of the angles of all triangles equals the sum of the interior angles of the $n$-gon, i.e. $(n-2)\cdot 180^\circ$.

    \Image{angles-polygon-fan.pdf}

    \Remark Proof~2 is an \uv{unrolled} version of Proof~1: induction cuts off triangles one by one, and here we see them all at once.

    \NamedProof{Proof 3 (interior point).} Choose a point $O$ in the interior of the $n$-gon and connect it to every vertex. This produces $n$ triangles with a total angle sum of $n\cdot 180^\circ$. At each vertex $A_k$ of the $n$-gon, the triangle angles at $A_k$ together form the entire interior angle of the $n$-gon there; the triangle angles at~$O$ together form a full angle of $360^\circ$. So the sum of the interior angles of the $n$-gon is
    $$
    n\cdot 180^\circ - 360^\circ = (n-2)\cdot 180^\circ.
    $$

    \Image{angles-polygon-interior.pdf}

    \Remark{non-convex polygons} The formula $(n-2)\cdot 180^\circ$ also holds for non-convex polygons, provided we measure all interior angles toward the inside -- so we count angles greater than $180^\circ$ as such.

    A fan triangulation from a single vertex may fail -- some diagonals can pass outside the polygon. It is still true, however, that every simple $n$-gon can be split into $n-2$ triangles (not necessarily as a fan), which rescues the proof. The reason is induction via a triangle formed by some vertex $V$ and its neighbors $U$, $W$ that lies entirely inside our $n$-gon. Cutting $V$ off along the diagonal $UW$ gives an $(n-1)$-gon; iterating yields $n-2$ triangles. One can prove that such a triangle always exists (in fact, two of them); see the \Link[https://en.wikipedia.org/wiki/Two_ears_theorem]{two ears theorem}.

    As for Proof~3: it works only if there is an interior point $O$ visible from every vertex (i.e., the segments $OA_i$ lie entirely inside). For some non-convex polygons no such $O$ exists (try drawing one) -- in that case Proof~3 cannot be salvaged directly and one must fall back on the triangulation described in the previous paragraph.
}

Now let's state two more simple but useful auxiliary claims:

\Exercise{}{
    Prove that the sum of the angles in the figure equals $180^\circ$.

    \Image{angles-complementary-parallel-statement.pdf}
}{
    Let $\alpha$ denote the angle that the transversal $AB$ makes with the parallel lines on the side of $B$ (as in the figure). The corresponding angles $\alpha$ at~$A$ and $\alpha$ at~$B$ are equal. At vertex $A$ the angles $\alpha$ and $\beta$ form a linear pair, so
    $$
    \beta + \alpha = 180^\circ.
    $$

    \Image{angles-complementary-parallel-solution.pdf}
}

\Exercise{}{
    Prove that the angle marked with a question mark in the figure equals $\alpha+\beta$.

    \Image{angles-complementary-in-triangle.pdf}
}{
    From the angle sum in the triangle, $\angle ACB = 180^\circ - \alpha - \beta$. The exterior angle at $C$ is in a linear pair with $\angle ACB$:
    $$
    180^\circ - \angle ACB = \alpha + \beta.
    $$

    \Image{angles-complementary-in-triangle.pdf}
}

The last two exercises may have looked very trivial. The reality, however, is that in practical problems they are very useful -- when we have a huge non-trivial diagram and have to chain together a large number of angle operations, it is genuinely advantageous to simplify even a small piece of that work. In the parallel-lines case we don't have to introduce an auxiliary corresponding angle, and in the triangle case we don't have to work with~$180^\circ$.

\sec Basics of the World of Lengths

So far we have only operated in the world of angles and have not dealt with lengths at all. Geometry, however, becomes interesting once we start connecting these worlds. Our foundation will be \textit{triangle congruence}. Let us recall the criteria.

\Highlight{
    Triangle congruence theorems:
    \begitems \style i
    \i $SSS$: two triangles are congruent if they agree on all three sides.
    \i $SAS$: two triangles are congruent if they agree on two sides and the angle between them.
    \i $ASA$: two triangles are congruent if they agree on one side and any two angles.
    \i $SsA$: two triangles are congruent if they agree on two sides and the angle opposite the longer of these sides.
    \enditems
}

Let us emphasize that in $SAS$ it really is important that the equal angle be the one between the two sides -- without this assumption, congruence may fail; see the figure. There we have two triangles $ABC$ and $A'B'C'$ for which $BC = B'C'$, $AC = A'C'$, and $\beta = \beta'$, yet they are clearly not congruent.

\Image{angles-sas-warning.pdf}

The $SsA$ theorem would be our rescue -- but here we cannot apply it, because the side opposite the angle $\beta$, namely $AC$, is not the longest, since it is clearly shorter than $BC$ \SmileWithTear.

Before we go further, let us understand what these theorems mean. In my view, a good way to look at congruence is the following: we want to construct a triangle when given some three of its elements -- will all the constructible triangles be congruent?

\Example{}{
    Convince yourself of the validity of the $SSS$ congruence theorem.
}{
    Imagine we are given three segments of lengths $a$, $b$, $c$ from which a triangle can be built. Without loss of generality, assume that $a=\max\{a,b,c\}$. Begin constructing the triangle from the segment~$BC$ of length~$a$. Then construct two circles: (a) the circle centered at~$B$ with radius $c$; (b) the circle centered at~$C$ with radius $b$. Thanks to $a \ge b$ and $a \ge c$, both of these circles cross the segment~$BC$. If $a=b+c$, they meet at exactly one point on it (which does not correspond to a triangle); if $a>b+c$, they fail to meet at all; and if $a<b+c$, they meet at two points $A$ and $A'$.

    \Image{angles-sss-proof.pdf}

    Finally, observe that the triangles $ABC$ and $A'BC$ are clearly mirror images of each other across $BC$. Likewise, had we started from a different segment, we would only produce a translation/rotation of this configuration.

    \Remark Note that along the way we essentially proved the triangle inequality. It is usually formulated as: the sum of any two sides of a triangle is greater than the third. We proved that the sum of the two shortest is greater than the longest. That clearly suffices, since the sum of the longest and any other is greater than the remaining one.
}

We can justify the other criteria similarly:

\Exercise{}{
    Convince yourself of the validity of the $SAS$ congruence theorem.
}{
    Suppose the sides $b = AC$, $c = AB$ and the angle $\angle BAC = \alpha$ between them are given. Construct vertex~$A$ and the segment $AB$ of length~$c$ -- this is determined up to congruence. Point $C$ must lie on a ray from~$A$ making angle~$\alpha$ with $AB$ and at distance~$b$ from~$A$. There are two rays from~$A$ making angle $\alpha$ with $AB$ (one on each side of $AB$): on each there is exactly one point at distance $b$ from~$A$, giving points $C$ and $C'$. The triangles $ABC$ and $ABC'$ are mirror images across $AB$, hence congruent.

    \Image{angles-sas-proof.pdf}
}

\Exercise{}{
    Convince yourself of the validity of the $ASA$ congruence theorem.
}{
    First, note that it doesn't matter which two angles are equal -- if two triangles agree on two angles, then the third is determined uniquely, since all three angles sum to $180^\circ$.

    For the purposes of our proof, suppose the equal angles are exactly the ones adjacent to the equal side. Let them be the angle $\beta$ at vertex $B$, $\gamma$ at vertex $C$, and the side $a = BC$ between them. Construct the segment $BC$ of length~$a$ -- this is determined up to congruence. Point $A$ must lie on the ray from~$B$ making angle $\beta$ with $BC$ and on the ray from~$C$ making angle $\gamma$ with $CB$. The sum of the interior angles at $B$ and $C$ in a triangle is less than $180^\circ$, i.e. $\beta + \gamma < 180^\circ$, so the two rays are not parallel -- they meet at a single point~$A$. The triangle $ABC$ is therefore determined uniquely (up to congruences).

    \Image{angles-asa-proof.pdf}
}

The less well-known $SsA$ theorem is a touch harder, so we will provide hints \SlightSmile.

\Problem{0}{$SsA$ congruence}{
    Convince yourself of the validity of the $SsA$ congruence theorem. In addition, figure out where the construction would break down if the angle were not opposite the longer of the two sides -- what would happen if both sides were equal?
}{
    Begin the construction with the shorter segment. How does the construction proceed once it is built?
}{
    After constructing the shorter segment, draw a ray at our given equal angle. One last step of the construction remains. Think about when we get no triangle, when one, and when two satisfying triangles.
}{
    Suppose the sides $a = BC$, $b = AC$ and the angle $\angle ABC = \beta$ opposite the longer side $b$ are given. We begin by constructing the shorter segment $BC$ of length~$a$. Point $A$ must lie on the ray from~$B$ making angle $\beta$ with $BC$ and at distance $b$ from~$C$, i.e. on the circle $k$ centered at~$C$ with radius $b$.

    The distance from~$B$ to the center of the circle is $BC = a$. Since $a < b$, the point~$B$ lies \textit{inside} the circle $k$. The ray from~$B$ thus starts inside the circle and continues outward, so it crosses the circle exactly once. Point $A$ is therefore determined uniquely, and the triangle $ABC$ is also unique (up to congruence).

    \Image{angles-Ssa-proof-one-solution.pdf}

    Let's see why the condition $b > a$ is essential. If $a = b$, the point~$B$ would lie directly \textit{on} the circle~$k$. The ray from~$B$ would meet it at $B$ itself (which does not give a triangle) and at exactly one further point -- which is the desired~$A$. The construction therefore still leads to a unique (isosceles) triangle, so the $SsA$ theorem formally holds even in this borderline case. We do not, however, list it under $SsA$ as a separate case: if we know that the triangle is isosceles with $BC=AC$ and we know one angle $\beta$, the other two angles are already determined by the isosceles property (both base angles equal~$\beta$). The same configuration is therefore covered by $SAS$ or $SSS$, and $SsA$ adds no new information.

    \Image{angles-Ssa-proof-isosceles.pdf}

    If $b < a$, the point $B$ would lie \textit{outside} the circle $k$. The ray could meet it at two points -- giving two different (non-congruent) triangles.

    \Image{angles-Ssa-proof-two-solutions.pdf}

    \Remark The proof can also be done by first constructing the longer segment. That, however, requires knowing the locus of points subtending a fixed angle from a fixed segment -- we will get to this in another handout.
}

Throughout our geometric journey we will encounter all of these statements. For now, however, let us show some concrete applications. We start with the most obvious statement, which still deserves to be proved:

\Theorem{isosceles triangle}{
    Prove that if a triangle~$ABC$ satisfies $AB=AC$, then $\angle ABC=\angle ACB$. Prove also the converse implication (i.e., that equality of angles implies equality of lengths).
}{
    \textit{Forward implication.} Assume $AB = AC$. Let $M$ be the midpoint of $BC$. The triangles $ABM$ and $ACM$ are congruent by $SSS$ ($AB=AC$ by assumption, $BM=CM$ by definition of the midpoint, common side $AM$), from which $\angle ABC = \angle ACB$. This proof also works with other choices of $M$. If we take $M$ to be the foot of the perpendicular from~$A$ to $BC$, we obtain the common side $AM$, $AB=AC$ and a right angle at~$M$, which lies opposite the longest side (the hypotenuse) $AB$, resp.~$AC$ -- we apply $SsA$. If we take $M$ to be the foot of the bisector of the interior angle at~$A$, we have $AB=AC$, the equal angle $\angle BAM=\angle CAM$, and the common side $AM$ between them -- we apply $SAS$.

    \Image{angles-isosceles-triangle.pdf}

    \textit{Converse implication.} Assume $\angle ABC = \angle ACB = \beta$. Let $M$ be the foot of the perpendicular from~$A$ to $BC$. By $ASA$, $\triangle ABM \cong \triangle ACM$, since we have two equal angles $\angle ABM = \angle ACM$ and $\angle AMB = \angle AMC = 90^\circ$ and the common side~$AM$, so $AB = AC$. Similarly, we could take $M$ to be the foot of the bisector of the interior angle at~$A$ -- from the equal angles $\angle BAM=\angle CAM$ and $\angle ABM=\angle ACM$ and the common side $AM$ we again apply $ASA$. Taking $M$ to be the midpoint of $BC$, however, doesn't help here: we get $BM=CM$, $AM$ common, and $\angle ABM=\angle ACM=\beta$, but the angle $\beta$ lies opposite $AM$, which may be shorter than $BM$ (when the angle at~$A$ is obtuse), so even $SsA$ cannot be applied in general -- the lesson is that even with the right point, its precise definition matters (we will see this many more times).

    \Remark Another cute proof is based on showing that $\triangle ABC \cong \triangle ACB$ (note the different order of vertices). In the case where we know the equal sides, we use $SAS$ or even $SSS$ -- which then gives equal angles. In the case where we know the equal angles, it will instead be $ASA$, and the congruence gives equal sides.
}

Try proving these simple facts rigorously -- it will be easier this time:

\Exercise{}{
    An equilateral triangle has three equal angles of $60^\circ$.
}{
    In $\triangle ABC$ with $AB = AC = BC$, the isosceles triangle theorem gives $\angle ABC = \angle ACB$ (from $AB=AC$); analogously, from $AB=BC$ we get $\angle BAC = \angle BCA$. All three angles are therefore equal, and from the angle sum of $180^\circ$ each is $60^\circ$.

    \Image{angles-equilateral.pdf}
}

\Exercise{}{
    An isosceles right triangle has angles $90^\circ$, $45^\circ$, $45^\circ$.
}{
    Suppose the right angle is at vertex~$A$, i.e. $\angle BAC = 90^\circ$, and the legs are equal, $AB = AC$. By the isosceles triangle theorem, $\angle ABC = \angle ACB$. Denote this common measure $\beta$. From the angle sum in the triangle,
    $$
    90^\circ + \beta + \beta = 180^\circ,
    $$
    so $2\beta = 90^\circ$, i.e. $\beta = 45^\circ$.

    \Image{angles-isosceles-right-triangle.pdf}
}

A very handy fact useful in non-trivial problems is the following. One possible proof goes via trigonometry. But we want a nice geometric one.

\Problem{0}{}{
    A right triangle has its other two angles equal to $60^\circ$ and $30^\circ$ if and only if its hypotenuse is twice as long as the shorter leg.
}{
    Say $AB$ is the hypotenuse and is twice as long as the leg $BC$. The trick is to introduce a point $B'$ such that $C$ is the midpoint of segment $BB'$.
}{
    Let our triangle have its right angle at vertex~$C$.

    \textit{($\Rightarrow$)} Assume $AB = 2\cdot BC$. Let $B'$ be the point such that $C$ is the midpoint of segment $BB'$. Then $B'C = BC$ and $C$ lies between $B$ and $B'$, so $\angle B'CA$ is in a linear pair with $\angle BCA = 90^\circ$, hence also equals $90^\circ$. We have $\triangle CBA \cong \triangle CB'A$ by $SAS$, since the angles at~$C$ are both right, the side~$AC$ is common, and $BC=B'C$.

    \Image{angles-90-60-30.pdf}

    Since $C$ is the midpoint of $BB'$, we have $BB' = 2\cdot BC$, and from the assumption $AB = 2\cdot BC$. Together,
    $$
    BB' = AB = AB',
    $$
    so $\triangle ABB'$ is equilateral and all its angles are $60^\circ$. In particular, $\angle ABC = \angle ABB' = 60^\circ$. From the angle sum in $\triangle ABC$ we compute $\angle BAC = 180^\circ - 90^\circ - 60^\circ = 30^\circ$.

    \textit{($\Leftarrow$)} It is not hard to see that the reasoning of the previous paragraph can easily be reversed -- this time the key congruence $\triangle CBA \cong \triangle CB'A$ comes from $ASA$.

    \Remark Another solution is to note that by Thales' theorem, the circumcenter~$O$ of triangle $ABC$ is also the midpoint of the hypotenuse $AB$. We will return to this solution in more detail when we discuss circles.
}

\sec Bisectors

In school we have surely encountered two kinds of bisectors: the perpendicular bisector of a segment and the bisector of an angle. In this section we will clarify familiar notions and important properties that we will use throughout.

\secc Perpendicular Bisector

We can define the perpendicular bisector of a segment as the line perpendicular to the given segment that passes through its midpoint. Another definition is that it is the set of points that are equidistant from the endpoints of our segment. But are these definitions equivalent? Well, it is not evident, so let us justify it.

\Theorem{perpendicular bisector}{
    The set of points equidistant from the endpoints of a segment $AB$ is the line perpendicular to $AB$ passing through its midpoint~$M$.
}{
    \textit{($\Rightarrow$)} Let point~$X$ lie on the perpendicular to $AB$ through the midpoint~$M$. The triangles $AMX$ and $BMX$ are congruent by $SAS$: $MA=MB$ (midpoint), common side $MX$, and equal right angles at~$M$. Hence $XA=XB$.

    \Image{angles-perpendicular-bisector-forward.pdf}

    \textit{($\Leftarrow$)} Suppose point~$X$ satisfies $XA=XB$ but does not lie on the perpendicular to $AB$ through~$M$. If $X$ lies on the line $AB$, it is clear that the only point of this line equidistant from~$A$ and~$B$ is precisely the midpoint~$M$.

    Suppose then that $X$ does not lie on the line $AB$. Without loss of generality, let our perpendicular be vertical and let~$X$ lie to its left. Denote by $Y$ the intersection of segment $XB$ with this perpendicular. By the already proved implication, $YA=YB$.

    \Image{angles-perpendicular-bisector-converse.pdf}

    In the isosceles $\triangle XAB$ ($XA=XB$) we have $\angle XAB=\angle XBA$. Since $Y$ lies on segment $XB$, the rays $BX$ and $BY$ coincide, so $\angle XBA=\angle YBA$. But in the isosceles $\triangle YAB$ ($YA=YB$) we have $\angle YBA=\angle YAB$. Altogether,
    $$
    \angle XAB=\angle XBA=\angle YBA=\angle YAB.
    $$
    This, however, is nonsense: point~$Y$ lies inside segment $XB$, so the ray $AY$ points into the interior of triangle $XAB$ and lies strictly inside the angle $\angle XAB$, i.e. $\angle YAB<\angle XAB$.
}

The most fundamental theorem involving perpendicular bisectors of the sides is, of course, the following statement:

\Theorem{existence of~$O$}{
    The perpendicular bisectors of the sides of any triangle meet at a single point. This point is the circumcenter of the triangle.
}{
    Let $ABC$ be a triangle. The perpendicular bisector of segment $AB$ is perpendicular to $AB$, and the perpendicular bisector of segment $AC$ is perpendicular to $AC$; these two perpendiculars are parallel only when $AB \parallel AC$, which clearly does not hold. Hence the perpendicular bisectors of segments $AB$ and $AC$ meet at a single point -- call it~$O$.

    By the previous theorem, $OA=OB$ (since $O$ lies on the perpendicular bisector of $AB$) and $OA=OC$ (since $O$ lies on the perpendicular bisector of $AC$). Together this gives $OB=OC$, which in turn tells us that $O$ also lies on the perpendicular bisector of segment $BC$. All three perpendicular bisectors therefore pass through~$O$.

    \Image{angles-circumcenter.pdf}

    \Remark{1} The point $O$ satisfies $OA=OB=OC$, so it is the center of a circle passing through all three vertices -- that is, the circumscribed circle (circumcircle) of triangle $ABC$.

    \Remark{2} The notation $O$ for the circumcenter is the international standard. In Czechia and Slovakia, the notation $S$ is also used (from the word \uv{střed}, meaning \uv{center}).
}

\secc Angle Bisector

As with the perpendicular bisector, let us first reflect on the definition. Take an angle $XAY$. By its \textit{bisector} (more precisely the \textit{internal bisector}) we mean a ray~$AZ$ such that $\angle XAZ=\angle ZAY$. The point~$Z$ is thus chosen so that the ray~$AZ$ splits the angle $XAY$ into two equal angles. Loosely, the entire line determined by this ray is also called the angle bisector.

\Image{angles-bisector-definition.pdf}

Just as with the perpendicular bisector, this definition is often confused with the claim that the angle bisector is the set of interior points of the angle that are equidistant from its two sides. Here too this deserves a proof.

\Theorem{angle bisector}{
    A point~$Z$ inside the angle $XAY$ is equidistant from the lines $AX$ and $AY$ if and only if $\angle XAZ=\angle ZAY$.
}{
    \textit{($\Rightarrow$)} Let the ray~$AZ$ split the angle $XAY$ into two equal angles. Let $X'$, $Y'$ be the feet of the perpendiculars dropped from~$Z$ to the lines $AX$, $AY$. The right triangles $AX'Z$ and $AY'Z$ are congruent by~$ASA$: common hypotenuse $AZ$, equal angles at~$A$, and equal right angles at~$X'$, $Y'$. Hence $ZX'=ZY'$, which are precisely the distances from~$Z$ to the lines $AX$ and $AY$.

    \Image{angles-bisector-equidistant.pdf}

    \textit{($\Leftarrow$)} Conversely, suppose $ZX'=ZY'$, where $X'$, $Y'$ are again the feet of the perpendiculars from~$Z$ to the lines $AX$, $AY$. The right triangles $AX'Z$ and $AY'Z$ are congruent by~$SsA$: common hypotenuse $AZ$, equal legs $ZX'=ZY'$, and equal right angles at~$X'$, $Y'$ opposite this hypotenuse (the longest side). Hence $\angle X'AZ=\angle Y'AZ$, which is precisely $\angle XAZ=\angle ZAY$.
}

Before we fully enjoy this theorem, let us also introduce the notion of the \textit{external bisector of an angle}. By the external bisector of an angle $XAY$ we mean the union of the bisectors of the two angles adjacent (in a linear pair) to~$XAY$ -- call them $XAY'$ and $X'AY$, where $X'$, $Y'$ lie on the rays opposite to $AX$, $AY$. The bisectors of these two adjacent angles are opposite rays issuing from~$A$; their union is therefore a line through~$A$.

\Image{angles-bisector-external-definition.pdf}

Before going further, let us prove this basic property:

\Exercise{}{
    Prove that the internal bisector of an angle is perpendicular to its external bisector.
}{
    Let $\beta$ denote half the angle $XAY$ and $\alpha$ half the angle adjacent to~$XAY$, as in the figure. The sum of all four angles along the line~$AY$ is
    $$
    \beta+\beta+\alpha+\alpha=180^\circ,
    $$
    from which $\alpha+\beta=90^\circ$. But $\alpha+\beta$ is precisely the angle between the internal and external bisectors, so they are perpendicular.

    \Image{angles-bisectors-perpendicular.pdf}
}

The external bisectors naturally also have the property that their points are equidistant from both sides of the angle -- after all, they are still bisectors of some angles. So if we look for the set of all points equidistant from two non-parallel \textit{lines} (i.e., not merely rays), we obtain the \textit{union} of the internal and external bisectors of the angle these lines form.

\Image{angles-two-lines-bisectors.pdf}

Let us now turn to the most famous consequence concerning angle bisectors. It is about the triangle, and we must pay attention to which bisector is internal and which is external.

\Theorem{existence of~$I$}{
    The internal angle bisectors of any triangle meet at a single point. This point is the incenter of the triangle.
}{
    Let $ABC$ be a triangle. Let~$I$ be the intersection of the internal angle bisectors at vertices~$B$ and~$C$, which clearly exists and clearly lies inside $ABC$.

    By the previous theorem, $d(I, AB) = d(I, BC)$ (since $I$ lies on the bisector of the angle at~$B$) and $d(I, BC) = d(I, AC)$ (since $I$ lies on the bisector of the angle at~$C$). Hence $d(I, AB) = d(I, AC)$. Since~$I$ lies inside the angle~$BAC$, this means that $I$ lies on its bisector. All three internal bisectors therefore pass through~$I$.

    \Image{angles-incenter.pdf}

    \Remark{1} The point~$I$ is equidistant from all three sides, so it is the center of a circle tangent to all three sides from the inside -- the incircle of triangle $ABC$.

    \Remark{2} The notation~$I$ comes from the word \uv{incenter}. In the Czech-Slovak olympiad the letter~$S$ is often used instead (for \uv{střed}, meaning \uv{center}).
}

In problems, the \textit{excircles} appear very often as well.

\Theorem{existence of~$I_a$}{
    In a triangle $ABC$, the internal bisector of the angle at~$A$ and the external bisectors of the angles at~$B$, $C$ meet at a single point. This point is the center of the excircle of triangle $ABC$ opposite vertex~$A$.
}{
    For clarity of notation, denote by $X$, $Y$ points on the rays opposite to $CA$, $BA$. The external bisector of the angle at~$B$ is exactly the internal bisector of the angle $YBC$, and the external bisector at~$C$ is the internal bisector of the angle $XCB$. Let~$I_a$ be the intersection of these bisectors.

    From the angle bisector theorem applied to the angles $YBC$ and $XCB$ we obtain the equalities
    $$
    d(I_a, AB) = d(I_a, BC) = d(I_a, AC).
    $$
    The point~$I_a$ is therefore equidistant from the lines $AB$ and $AC$, so it lies on either the internal or the external bisector of the angle $BAC$. But it clearly lies inside the angle $BAC$, so it is on the internal bisector -- which therefore passes through~$I_a$.

    \Image{angles-excenter.pdf}

    \Remark{1} The point~$I_a$ is equidistant from the line $BC$ and from the lines containing the sides $AB$, $AC$, so it is the center of a circle tangent to the side $BC$ externally and to the other two sides (more precisely their extensions) internally -- the excircle of the triangle opposite vertex~$A$.

    \Remark{2} The notation~$I_a$ is the common international notation. The escribed circle is called the \uv{excircle}, so perhaps $E_a$ would have been more logical, but it is rarely used. Then again, the inscribed circle is associated with the incenter, so perhaps that's why \UpsideDownSmile.
}

Angle bisectors, both internal and external, are an extremely useful and beloved concept in competition problems. They often hide behind other conditions, and the key step in a solution is often to realize that some point lies on two bisectors, and therefore also on the third. Sometimes the bisectors are internal, sometimes external. The external ones are often harder to spot. Making friends with bisectors definitely pays off.

\sec What to Remember

\secc Techniques

\begitems \style N
\i With parallel lines, look for equal alternate and corresponding angles, or angles summing to~$180^\circ$.
\i It is useful to view the exterior angle of a triangle as the sum of the two non-adjacent interior angles.
\i To bridge the worlds of lengths and angles, congruence of two triangles is the right tool.
\i Auxiliary points (midpoint of a segment, mirror image, extension) often reveal a congruence or a special triangle.
\enditems

\secc Useful Facts

\begitems \style N
\i The angle sum in a triangle is $180^\circ$, in general in a convex $n$-gon $(n-2) \cdot 180^\circ$.
\i Triangle congruence theorems: $SSS$, $SAS$, $ASA$, $SsA$.
\i In a triangle, two sides are equal if and only if the angles opposite them are equal.
\i A right triangle has its hypotenuse twice as long as a leg if and only if its angles are $90^\circ, 60^\circ, 30^\circ$.
\enditems

\sec Problems

In this section you will find various problems for which the knowledge from this handout is sufficient -- no advanced topics like \textit{inscribed and central angles} are needed (we will look at those problems later).


\Problem{0}{MO district round Z8 2025}{
    In a triangle $ABC$, point $D$ lies on side $BC$ and point $E$ on side $AC$, with $AB = BE = EC = CD$ and $BD = DE$. Determine the measures of the angles $\angle ACB$ and $\angle BAD$.
}{
    Set $\gamma = \angle ACB$ and try to express every angle in the figure in terms of~$\gamma$. We have plenty of isosceles triangles.
}{
    One possible path is to express the angles $\angle CBE$, $\angle BED$, $\angle CDE$, $\angle CED$ in terms of $\gamma$ in turn. Can we now write down an equation somewhere from which $\gamma$ can be computed?
}{
    Look at the angle sum in triangle~$BEC$. It should let us compute $\gamma$.
}{
    Computing~$\angle BAD$ takes more steps. The first is to compute as many angles as possible and notice something.
}{
    Once we compute the angles in $ABC$, we see that $ABC$ is isosceles.
}{
    Now we can exploit symmetry; one option is to prove the congruence of the triangles $EAD$ and $DBE$.
}{
    \Answer{$\angle ACB = 36^\circ$ and $\angle BAD = 36^\circ$.}

    Set $\gamma = \angle ACB$.

    In the isosceles $\triangle BEC$ ($BE=EC$) we have $\angle EBC = \angle ECB = \gamma$. Since $D$ lies on $BC$, $\angle DBE = \angle EBC = \gamma$, and by the isosceles property of $\triangle BDE$ ($BD=DE$) also $\angle DEB = \gamma$. The exterior-angle theorem in $\triangle BDE$ at~$D$ gives
    $$
    \angle EDC = \angle DBE + \angle DEB = 2\gamma,
    $$
    and in the isosceles $\triangle CDE$ ($CD=CE$) we have $\angle CED = \angle CDE = 2\gamma$. From its angle sum $\gamma + 4\gamma = 180^\circ$, so $\gamma = 36^\circ$.

    \Image{angles-isosceles-chain.pdf}{0.8}

    Now note that $\angle AEB$ is in a linear pair with $\angle BEC = 3\gamma = 108^\circ$, so it equals $72^\circ$. By the isosceles property of $\triangle BAE$ we also get $\angle BAE = 72^\circ$. From $\triangle ABC$, where the angles at $C$ and $A$ are $36^\circ$ and $72^\circ$ respectively, we compute $\angle CBA = 72^\circ$, so triangle $ABC$ is isosceles.

    Now observe that from $CA = CB$ and $CE = CD$ we get $EA = DB$. Consider the triangles $EAD$ and $DBE$: they are isosceles with equal apex angle (equal to $180^\circ - 2\gamma$) and equal legs, so by $SAS$ they are congruent. This gives $\angle DAE = \angle DBE = 36^\circ$.

    Finally,
    $$
    \angle BAD = \angle BAE - \angle DAE = 2\gamma - \gamma = \gamma = 36^\circ.
    $$
}

\Problem{0}{DuoGeo 2025}{
    Inside the square $ABCD$, equilateral triangles $ABX$ and $CDY$ have been drawn. Determine the sum of the marked angles.

    \Image{angles-square-equilateral-statement.pdf}{0.7}
}{
    Try first to compute as many angles as possible in the figure.
}{
    The key to computing the final angle is to find a suitable isosceles triangle using equal lengths.
}{
    \Answer{$180^\circ$.}

    By symmetry, all four marked angles are equal. So it suffices to find the measure of one of them -- we focus on $\angle YDX$.

    First, $\angle XAD = \angle BAD - \angle BAX = 90^\circ - 60^\circ = 30^\circ$; analogously $\angle ADY = \angle ADC - \angle YDC = 90^\circ - 60^\circ = 30^\circ$.

    Next, $DA = AB = AX$ (the first is a side of the square, the second a side of the equilateral triangle $ABX$), so $\triangle AXD$ is isosceles with base $DX$. From $\angle XAD = 30^\circ$ and the angle sum of the triangle,
    $$
    \angle ADX = \tfrac{1}{2}\bigl(180^\circ - 30^\circ\bigr) = 75^\circ.
    $$
    The desired angle is the difference of the two already computed:
    $$
    \angle YDX = \angle ADX - \angle ADY = 75^\circ - 30^\circ = 45^\circ.
    $$
    The sum of all four marked angles is therefore $4 \cdot 45^\circ = 180^\circ$.

    \Image{angles-square-equilateral-solution.pdf}{0.8}
}

\Problem{0}{MO school round C 2024}{
    In a trapezoid $ABCD$ with $AB \parallel CD$, the interior angle bisectors at vertices $C$ and $D$ meet on segment~$AB$. Prove that $AD + BC = AB$.
}{
    Let~$P$ be the common intersection. We have parallel lines and angle bisectors, which gives plenty of equal angles.
}{
    The goal is to find isosceles triangles.
}{
    Let $P$ be the common intersection of the bisectors; by hypothesis it lies on segment~$AB$. We will show that
    $$
    AD = AP \quad \hbox{and} \quad BC = BP,
    $$
    from which $AD + BC = AP + BP = AB$ directly.

    Since $DP$ is the bisector of the interior angle at~$D$, we have $\angle ADP = \angle CDP$. By $AB \parallel CD$, $\angle CDP$ and $\angle APD$ are alternate angles cut by the transversal~$DP$, so
    $$
    \angle APD = \angle CDP = \angle ADP.
    $$
    In triangle $ADP$ the angles at vertices $D$ and $P$ are equal, so $AD = AP$. An analogous argument at vertex~$C$ yields $BC = BP$.

    \Image{angles-trapezoid-bisectors.pdf}{0.8}
}

\Problem{0}{MO school round A 2023}{
    In a convex pentagon $ABCDE$ we have $\angle CBA = \angle BAE = \angle AED$. On the sides $AB$ and $AE$ there exist points $P$ and $Q$ respectively such that $AP = BC = QE$ and $AQ = BP = DE$. Prove that $CD \parallel PQ$.
}{
    We have plenty of equal angles and sides; try to find congruent triangles.
}{
    The key congruences are among triangles $PBC$, $QAP$, and $DEQ$. They yield useful equal pieces. Keep looking for more congruences.
}{
    The final congruence to prove is between $CPQ$ and $DQP$. Why does that suffice?
}{
    Since $BC = AP = EQ$, $BP = AQ = ED$, and $\angle CBP = \angle PAQ = \angle QED$, the triangles $PBC$, $QAP$, and $DEQ$ are mutually congruent by $SAS$.

    \Image{angles-pentagon.pdf}{0.8}

    From this $CP = PQ = QD$, and also
    $$
    \gather
    \angle CPQ = 180^\circ - \angle BPC - \angle APQ = \\ = 180^\circ - \angle PQA - \angle EQD = \angle PQD.
    \endgather
    $$
    Hence, by $SAS$, the isosceles triangles $CPQ$ and $DQP$ are also congruent. From this, the altitudes from vertices $C$ and $D$ to the common side $PQ$ are equal; these altitudes are also parallel (both perpendicular to $PQ$), so $CD \parallel PQ$.
}

\Problem{0}{DuoGeo 2025}{
    A quadrilateral $ABCD$ is given, with $T$ the intersection of its diagonals. Assume that the angles $\angle BAC$ and $\angle DBA$ measure $30^\circ$ and $45^\circ$ respectively. On segment $BT$ there is a point $Z$ such that $CZ \perp BT$. Assume that line $CZ$ meets segment $AB$ at point $M$. Let $R$ be the intersection of segments $AT$ and $MD$. Suppose that $AM = AR$ and $MR + TD = 14$. Determine the length of segment $BZ$.
}{
    Keep computing angles until we find an isosceles triangle.
}{
    Patient angle chasing leads to $DTR$ being isosceles. That should shed light on the condition $MR + TD = 14$.
}{
    Once the isosceles relations are proved, the condition translates to $MD = 14$. That will help later -- for now we cannot compute more angles and must find something from the world of lengths. The key is to find a nice right triangle.
}{
    One can show that the angles of triangle $MDZ$ are $90$-$60$-$30$. We know its hypotenuse $MD$. Our nicely proved auxiliary statement now gives another length. From there it is one step to the solution.
}{
    \Answer{$BZ = 7$.}

    In triangle $ATB$ we know the angles at vertices $A$ and $B$, namely $30^\circ$ and $45^\circ$. Hence the exterior angle at~$T$ equals the sum of these angles; concretely
    $$
    \angle DTR = \angle TAB = \angle TBA = 30^\circ + 45^\circ = 75^\circ.
    $$
    In the isosceles triangle $AMR$ with $AM = AR$, the angle at vertex $A$ equals $30^\circ$, so both angles at the base $MR$ have measure $\tfrac{1}{2}(180^\circ - 30^\circ) = 75^\circ$. In particular $\angle ARM = 75^\circ$, and so its vertical angle $\angle DRT$ also has measure $75^\circ$.

    \Image{angles-quad-30-45.pdf}{0.8}

    Triangle $DTR$ therefore has two angles of measure $75^\circ$, so it is isosceles with base $TR$ and $DT = DR$. Hence
    $$
    MD = MR + RD = MR + TD = 14.
    $$
    Moreover, $\angle TDR = 180^\circ - 2 \cdot 75^\circ = 30^\circ$.

    Now consider triangle $MDZ$. Since $T$, $Z$ lie on segment $BD$ and $R$ lies on segment $MD$, we have $\angle MDZ = \angle RDT = 30^\circ$. Triangle $MDZ$ is therefore right-angled (at~$Z$) with hypotenuse $MD$ and an angle of $30^\circ$ at~$D$, so by the earlier statement about the $30^\circ$-$60^\circ$-$90^\circ$ triangle
    $$
    MZ = \tfrac{1}{2} MD = 7.
    $$

    Finally, triangle $MZB$ has a right angle at~$Z$ and an angle of $45^\circ$ at~$B$, so $\angle BMZ = 45^\circ$. It is therefore an isosceles right triangle and $BZ = MZ = 7$.
}

\DisplayExerciseSolutions

\DisplayHints

\DisplayProblemSolutions

\bye
